\section{引言}
\label{intro:section}

纠缠投影测量同时受到全局基态正交性和局部约化态几何的约束。优雅联合测量（elegant joint measurement，EJM）将这两种性质结合起来：局部 Bloch 向量位于正四面体的顶点，而对应的全局态构成正交归一基。$n=2$ 的 EJM 最初出现于量子网络研究~\cite{Gisin2017}，并属于一族连接它与 Bell 态测量的连续测量~\cite{Tavakoli2021}。改变四面体半径会改变单比特约化谱，进而改变每个单比特与其余系统划分下的纠缠。将这种调节推广到更多比特，需要所有测量结果的局部几何始终与全局正交性兼容。

Pauli 轨道构造为这一问题提供了适合的框架~\cite{PauwelsPauli2026,PauwelsEJM2026}。在相应的相位正规形中，每个对角酉算符都会生成完整正交归一基，剩余约束集中于局部 Bloch 向量。Pauwels 与 Gisin 对每个 $n\ge2$ 给出了显式基准 EJM，其共同半径为
\begin{equation}
 r_E(n)=\frac{\sqrt3}{2^{n-1}},
 \label{intro:reference-radius}
\end{equation}
并构造了在 $|t|\le\pi/2$ 上满足 $r_E(n)\cos t$ 的偶数比特相位族~\cite{PauwelsGisin2026}。该文已将局部信号的来源解释为涉及单个计算基相位修正的干涉。其奇数比特分析将 $n\ge5$ 时经过指定基准点的显式可调族留作开放问题，并通过数值 Jacobian 秩识别出$n=3$ 时的局部孤立基准点。其它对称多比特构造允许更宽的局部几何~\cite{Ding2025,He2026}；本文保持每个比特上的一个固定正四面体，并要求所有比特和结果具有共同半径。

本工作的上一版对每个奇数 $n\ge5$ 给出了半径为 $r_E(n)(1-\sin t)$ 的显式增强族~\cite{Zhang2026Odd}。本文的构造不再区分奇偶：一个依赖四个二进制特征的新 Boolean 核，使同一半径规律适用于全部整数 $n\ge5$。它在 $t=0$ 逐列恢复基准测量，并连续连接局部极化为零与基准半径两倍这两个端点。这一增强自然引出其极限问题。为此，我们考察所有经过基准测量的两值实相位生成元，分类它们所允许的半径规律，证明两倍基准半径在该模型内最优；这一模型也包括四特征构造之外的生成元。

统一构造还使新增强变形与已知偶数比特余弦变形能够在同一相位表示中比较。我们考察两种调节能否在保持局部正四面体的同时独立进行。对每个偶数 $n\ge8$，一个显式的完整测量二维族给出了肯定答案，其两个坐标截面恰为这两种变形。$n=6$ 时，混合干涉阻碍了这一扩张；在相应的四颜色相位模型内，局部解只有两条相交分支。

更低维的情形进一步说明了直线相位路径何时不再足够。在 $n=2,3,4$ 时，两值相位直线均不能增强基准半径。$n=3$ 时，我们将 Pauwels 与 Gisin 的数值证据加强为精确的局部刚性证书。$n=4$ 时，即使允许任意实生成元，所有允许的直线仍不能增强半径，但一条显式弯曲路径可将半径增加至 $(1+1/\sqrt2)r_E(4)$。因此，在同一相位表示中改变路径形状即可突破$n=4$ 时的直线障碍。

第~II~节建立将这些问题归约为局部干涉条件的相位框架。第~III、IV~节给出统一构造及 Boolean 模型内的精确界。第~V~节建立它与已知偶数比特族的关系，第~VI~节解释低维障碍并给出$n=4$ 时的弯曲构造。各附录给出完整证明。
